\section{Summary}
\label{sec:summary}

The system described here designs its own exploratory test cases for a mobile
or web application, executes them against a running build without supervision,
and writes what it learns into a knowledge graph that shapes the next test. No
human writes a test case, a selector or a script. It is the
successor to the static-knowledge test generator described in section 2 of
ETA-TR-2026-001, and it implements all eight of the enhancement requirements in
that document.

\subsection*{What the system does}

A campaign begins with whatever knowledge exists: a requirements
specification, a design export, a defect history, or nothing at all. Three
language-model roles then operate in a loop over one Neo4j graph. A
\emph{planner} decides what is worth testing next and writes an objective in
plain language. A \emph{player} drives the application toward that objective on
a real device or in a browser, recording every step. An \emph{investigator}
reads the resulting trajectory against what the test intended and reduces it to
durable, deduplicated findings. Each round leaves the graph richer than it found
it, so the next objective is chosen from better facts.

The property that distinguishes this system from the work surveyed in
\Cref{sec:positioning} is not any single mechanism. It is that a failed run is
classified, at the point of failure, as one of three things: the application
misbehaved, the agent failed, or the environment broke. Until that distinction
exists, a defect count cannot be interpreted, because ``the test failed'' and
``the agent got lost'' are the same observation.

\subsection*{What has been achieved}

\begin{table}[H]
  \centering\small
  \begin{tabularx}{\textwidth}{@{}>{\hsize=0.52\hsize}Y>{\hsize=0.38\hsize}Y>{\hsize=1.10\hsize}Y@{}}
    \toprule
    Property & Result & Basis \\
    \midrule
    Requirement families delivered & 8 of 8 & All ETA-REQ-301 to 308, verifiable on the dashboard \\
    \addlinespace
    Autonomy, best campaign & 100\,\% & 64 web test cases with nothing attributed to the agent; and separately 25 Android executions with none lost \\
    \addlinespace
    Autonomy, hardest target & 69\,\% & 45-test campaign against a production marketplace with no stable identifiers \\
    \addlinespace
    Cost per executed test & USD 0.008 & Metered billing over 88 device executions, not a price-card estimate \\
    \addlinespace
    Route efficiency & 50 steps to 9 & Median device steps per test, before and after navigation memory \\
    \addlinespace
    Targets driven & 5 & Two Android applications and three web applications, no code change between them \\
    \addlinespace
    Execution record & 123 trajectories & Every device execution preserved step by step and re-readable \\
    \bottomrule
  \end{tabularx}
\end{table}

\keypoint{The figure we would ask you to read first is not the pass rate. It is
autonomy: the share of executions that end in a statement about the application
rather than a failure of our agent. It moved from effectively zero to 100\,\%
over two campaigns, and it is the precondition for every other number in
this report meaning anything.}

\subsection*{What has not been achieved}

Precision and recall for defect detection are not reported, and cannot be. Both
require a build whose defects are known in advance. No defect export and no
seeded-defect build were available, so every application-level failure in this
report is a \emph{candidate} defect that a human must confirm. The system is
built to say so rather than to present candidates as findings.

This is not a gap peculiar to this project. Of the seven comparable systems
reviewed in \Cref{sec:positioning}, none reports precision or recall either.
\Cref{sec:next} sets out the one experiment that would close the question and
what we would need in order to run it.
